% بنك التمارين — الفصل 2
% يُحدَّد الفصل من اسم هذا الملف.

\begin{exo}[من الجليد البارد إلى الماء المغلي: رتب المقدار]{chauffage-glace-eau-ebullition}
\seoready{true}
\keywords{المسعرية, تغير الحالة, الانصهار, التبخر, الحرارة الكامنة, رتبة المقدار}

\textbf{المعطيات العددية (عند الضغط الجوي):}
\[
\begin{aligned}
c_{\text{ice}} &= 2{,}1\times10^3\,\text{J}\!\cdot\!\text{kg}^{-1}\!\cdot\!\text{K}^{-1},
& L_f &= 3{,}34\times10^5\,\text{J}\!\cdot\!\text{kg}^{-1},\\
c_{\text{water}} &= 4{,}18\times10^3\,\text{J}\!\cdot\!\text{kg}^{-1}\!\cdot\!\text{K}^{-1},
& L_v &= 2{,}26\times10^6\,\text{J}\!\cdot\!\text{kg}^{-1}.
\end{aligned}
\]
عند الضغط الجوي، يُسخَّن تدريجياً كيلوغرام واحد من الجليد، ودرجة حرارته الابتدائية $-100\,{}^\circ\text{C}$.
\begin{enumerate}
  \item احسب الطاقة اللازمة لتسخين الجليد من $-100\,{}^\circ\text{C}$ إلى $0\,{}^\circ\text{C}$.
  \item احسب الطاقة اللازمة لصهر الجليد كلياً عند $0\,{}^\circ\text{C}$.
  \item احسب الطاقة اللازمة لتسخين الماء السائل من $0\,{}^\circ\text{C}$ إلى $100\,{}^\circ\text{C}$.
  \item احسب الطاقة الإضافية اللازمة لتبخير الماء كلياً عند $100\,{}^\circ\text{C}$.
  \item أي مرحلة تتطلب أكبر طاقة؟ قارن طاقة تسخين كيلوغرام واحد من الماء من $0\,{}^\circ\text{C}$ إلى $100\,{}^\circ\text{C}$ بطاقة الوضع لكتلة $m$ تسقط من ارتفاع عشرة أمتار. ما الكتلة التي ينبغي إسقاطها لتحرير الطاقة نفسها؟
  \item ادرس الآن العملية العكسية. تتكاثف كتلة $m_v$ من بخار الماء على الزجاج الأمامي، من دون أن يبرد الماء السائل الناتج بعد ذلك. اكتب الحرارة $Q_{\text{rel}}$ المتحررة أثناء التكاثف واحسبها من أجل $m_v=1{,}0\times10^{-2}\,\text{kg}$. عند درجة حرارة الزجاج خذ $L_v=2{,}45\times10^6\,\text{J}\!\cdot\!\text{kg}^{-1}$.
\end{enumerate}
\begin{indication}
لتسخين طور من دون تغير الحالة استخدم $Q=mc\Delta T$، ولتغير الحالة عند درجة حرارة ثابتة استخدم $Q=mL$. طاقة الوضع الثقالية لكتلة $m$ على ارتفاع $h$ هي $E_p=mgh$؛ خذ $g=9{,}81\,\text{m}\!\cdot\!\text{s}^{-2}$.
\end{indication}
\begin{solution}
\textbf{1.} يتطلب تسخين الجليد إلى درجة انصهاره
\[
Q_1=mc_{\text{ice}}\Delta T
=1{,}0\times2{,}1\times10^3\times100=2{,}1\times10^5\,\text{J}.
\]
\textbf{2.} تبقى درجة الحرارة $0\,{}^\circ\text{C}$ أثناء الانصهار:
\[
Q_2=mL_f=1{,}0\times3{,}34\times10^5=3{,}34\times10^5\,\text{J}.
\]
\textbf{3.}
\[
Q_3=mc_{\text{water}}\Delta T
=1{,}0\times4{,}18\times10^3\times100=4{,}18\times10^5\,\text{J}.
\]
\textbf{4.} يتطلب التبخر الكامل أيضاً
\[
Q_4=mL_v=1{,}0\times2{,}26\times10^6=2{,}26\times10^6\,\text{J}.
\]
إذن تتطلب هذه المرحلة وحدها طاقة من رتبة بضعة ميغاجول.

\textbf{5.} التبخر هو، بفارق كبير، المرحلة الأكثر استهلاكاً للطاقة:
\[
Q_4=2{,}26\times10^6>Q_3=4{,}18\times10^5>Q_2=3{,}34\times10^5>Q_1=2{,}1\times10^5\,\text{J}.
\]
ويتطلب نحو $(2{,}26\times10^6)/(4{,}18\times10^5)\approx5{,}4$ أضعاف طاقة تسخين الماء السائل من $0$ إلى $100\,{}^\circ\text{C}$.

لكتلة $m$ تسقط من $h=10\,\text{m}$، لدينا $E_p=mgh$. بوضع $mgh=Q_3$ نحصل على
\[
m=\frac{Q_3}{gh}=\frac{4{,}18\times10^5}{9{,}81\times10}
\approx4{,}26\times10^3\,\text{kg}.
\]
يلزم إذاً إسقاط كتلة تقارب $4{,}3$ أطنان من عشرة أمتار لتحرير طاقة مساوية لتسخين كيلوغرام من الماء من $0$ إلى $100\,{}^\circ\text{C}$. هذه قيمة هائلة تفسر أهمية تسخين الماء في استهلاك الطاقة المنزلي. كما أن السعة الحرارية النوعية للماء كبيرة مقارنة بالمعادن الشائعة؛ فهي مثلاً نحو عشرة أمثال سعة النحاس (انظر التمارين التالية).

\textbf{6.} التكاثف هو العملية العكسية للتبخر، فيمنح البخار الزجاج الحرارة الكامنة:
\[
Q_{\text{rel}}=m_vL_v.
\]
من أجل $m_v=1{,}0\times10^{-2}\,\text{kg}$،
\[
Q_{\text{rel}}=1{,}0\times10^{-2}\times2{,}45\times10^6
=2{,}45\times10^4\,\text{J}.
\]
لذلك قد يحرر تكاثف كتلة صغيرة من البخار كمية غير مهملة من الحرارة يمتصها الزجاج ومحيطه.
\end{solution}
\end{exo}

\begin{exo}[نتح وتبخر شجرة كبيرة: مكيّف طبيعي؟]{odg-evapotranspiration-arbre}
\seoready{true}
\keywords{رتبة المقدار, النتح والتبخر, شجرة, الحرارة الكامنة, قدرة التبريد, مكيّف, COP}

في يوم حار وجاف، قد تفقد شجرة كبيرة تتلقى ماءً كافياً نحو $500\,\text{L}=5{,}0\times10^{-1}\,\text{m}^3$ من الماء بالنتح والتبخر. ولأن النتح يحدث أساساً في الضوء، نفترض أنه يستمر اثنتي عشرة ساعة. يغطي تاج الشجرة مساحة أرضية $A_{\text{tree}}=50\,\text{m}^2$. كثافة الماء وحرارة تبخره الكامنة عند الضغط الجوي هما
\[
\rho_{\text{water}}=1{,}0\times10^3\,\text{kg}\!\cdot\!\text{m}^{-3},
\qquad L_v=2{,}45\times10^6\,\text{J}\!\cdot\!\text{kg}^{-1}.
\]
\begin{enumerate}
  \item احسب كتلة الماء التي تفقدها الشجرة بالنتح والتبخر خلال اليوم والطاقة اللازمة لتبخيرها.
  \item فسّر لماذا تنتج هذه العملية تأثيراً مبرداً.
  \item احسب قدرة التبريد المتوسطة للشجرة لكل متر مربع خلال ساعات النتح النشط الاثنتي عشرة.
  \item للمقارنة، ادرس مكيّفاً بقدرة كهربائية $P_{\mathrm{el}}=2{,}0\times10^3\,\text{W}$ ومعامل أداء تبريد $\mathrm{COP}=3{,}0$، يبرد غرفة مساحتها $A_{\text{room}}=20\,\text{m}^2$. بالتعريف $\mathrm{COP}=P_{\text{cool}}/P_{\mathrm{el}}$. احسب قدرة التبريد والقدرة لكل متر مربع وقارنهما بالشجرة.
  \item لماذا لا يمكن استنتاج أن الشجرة والمكيّف يعطيان الإحساس نفسه بالتبريد، حتى إذا كانت قدرتاهما لكل وحدة مساحة من رتبة المقدار نفسها؟
  \item هل يمكن تبريد المنزل بكثير من النباتات الداخلية؟ (ثقافة عامة: يعتمد الجواب على عمليات حيوية.)
\end{enumerate}
\begin{indication}
استخدم $m=\rho V$ و$Q_{\mathrm{evap}}=mL_v$ و$P=Q/\tau$. للمكيّف استخدم $P_{\text{cool}}=\mathrm{COP}\,P_{\mathrm{el}}$.
\end{indication}
\begin{solution}
\textbf{1.}
\[
m=\rho_{\text{water}}V=1{,}0\times10^3\times5{,}0\times10^{-1}
=5{,}0\times10^2\,\text{kg},
\]
\[
Q_{\mathrm{evap}}=mL_v=5{,}0\times10^2\times2{,}45\times10^6
=1{,}225\times10^9\,\text{J}.
\]
أي إن رتبة المقدار غيغاجول واحد.

\textbf{2.} التبخر تغير طور ماص للحرارة؛ يأخذ الطاقة من الأوراق والتربة والهواء المحيط، فتقل درجة حرارة الأوراق ومحيطها القريب.

\textbf{3.} اثنتا عشرة ساعة تساوي $\tau=12\times3600=4{,}32\times10^4\,\text{s}$. لذلك
\[
P_{\text{tree}}=\frac{Q_{\mathrm{evap}}}{\tau}
=\frac{1{,}225\times10^9}{4{,}32\times10^4}\approx2{,}84\times10^4\,\text{W},
\]
ولكل وحدة مساحة
\[
\frac{P_{\text{tree}}}{A_{\text{tree}}}
=\frac{2{,}84\times10^4}{50}\approx5{,}7\times10^2\,\text{W}\!\cdot\!\text{m}^{-2}.
\]
\textbf{4.}
\[
P_{\text{cool}}=\mathrm{COP}\,P_{\mathrm{el}}
=3{,}0\times2{,}0\times10^3=6{,}0\times10^3\,\text{W},
\]
\[
\frac{P_{\text{cool}}}{A_{\text{room}}}
=\frac{6{,}0\times10^3}{20}=3{,}0\times10^2\,\text{W}\!\cdot\!\text{m}^{-2}.
\]
في هذا النموذج، تبلغ قدرة النتح والتبخر للشجرة لكل مساحة قرابة ضعفي قدرة مكيّف عادي.

\textbf{5.} تتعلق المقارنة بتدفقات الطاقة وحدها. يبرد المكيّف حيزاً مغلقاً ومضبوطاً، بينما ينتشر بخار الشجرة في الجو وتجلب الرياح هواءً دافئاً. يعتمد الإحساس على التهوية والرطوبة ودرجة الحرارة والمسافة من الشجرة. ويتغير النتح مع الشمس والرياح والرطوبة والنوع والماء المتاح، وقد ينخفض كثيراً في الجفاف. كما توفر الشجرة الظل فتقلل مباشرة الإشعاع الشمسي الواصل إلى الأرض والناس. لذا تقارن الحسابات رتب المقدار ولا تثبت تكافؤ الراحة الحرارية.

\textbf{6.} عملياً يكون تبريد النباتات الداخلية ضعيفاً جداً. لدى معظم النباتات يشجع الضوء فتح الثغور، وهي مسام صغيرة في الأوراق يخرج منها بخار الماء. تكون إضاءة الداخل ضعيفة عادة، فتفتح الثغور أقل وينخفض النتح. وفي غرفة قليلة التهوية تؤدي زيادة الرطوبة أيضاً إلى إبطاء التبخر تدريجياً.

قد تعطي نباتات كثيرة تبريداً تبخرياً محلياً طفيفاً، لكنها لا تعوض المكيّف. ولإخراج الطاقة باستمرار من المنزل يجب طرد الهواء الرطب إلى الخارج. قد يحفز الضوء الاصطناعي القوي النتح، لكن طاقته الكهربائية تنتهي تقريباً كلها حرارة في الغرفة. نباتات CAM المتكيفة مع الجفاف، ومنها الصبار، استثناء إذ تفتح ثغورها أساساً ليلاً.
\end{solution}
\end{exo}

\begin{exo}[مزج كتلتين من الماء]{calorimetrie-melange-eau}
\seoready{true}
\keywords{المسعرية, المزج, السعة الحرارية, الحرارة النوعية, جوزيف بلاك, درجة حرارة التوازن}

يُصب $m_1=0{,}200\,\text{kg}$ من الماء عند $T_1=80\,^\circ\text{C}$ في وعاء يحوي $m_2=0{,}300\,\text{kg}$ من الماء عند $T_2=15\,^\circ\text{C}$. الوعاء مسعر مثالي: نهمل فقد الحرارة إلى الخارج وسعته الحرارية. نذكر أن $Q=mc\Delta T$، حيث $c=4{,}18\times10^3\,\text{J}\!\cdot\!\text{kg}^{-1}\!\cdot\!\text{K}^{-1}$ هي الحرارة النوعية للماء السائل.
\begin{enumerate}
  \item بيّن أن الحرارة التي يفقدها الماء الساخن يستقبلها الماء البارد كاملة، واكتب معادلة الحفظ التي تضم $m_1$ و$m_2$ و$c$ و$T_1$ و$T_2$ ودرجة التوازن $T_{eq}$.
  \item استنتج عبارة $T_{eq}$ واحسب قيمتها.
  \item احسب الحرارة $Q$ التي يفقدها الماء الساخن أثناء المزج.
  \item هل يتيح مزج مادتين من النوع نفسه تحديد $c$؟ علّل.
\end{enumerate}
\begin{indication}
المسعر معزول وصلب، لذلك تحفظ الطاقة الداخلية الكلية $U_1+U_2$: ما يفقده أحد الجزأين يستقبله الآخر بإشارة معاكسة.
\end{indication}
\begin{solution}
\textbf{1.}
\[
m_1c(T_{eq}-T_1)+m_2c(T_{eq}-T_2)=0.
\]
\textbf{2.} يختزل العامل $c$:
\[
T_{eq}=\frac{m_1T_1+m_2T_2}{m_1+m_2}
=\frac{0{,}200\times80+0{,}300\times15}{0{,}200+0{,}300}=41\,^\circ\text{C}.
\]
\textbf{3.}
\[
Q=m_1c(T_1-T_{eq})=0{,}200\times4{,}18\times10^3\times(80-41)
\approx3{,}26\times10^4\,\text{J}.
\]
ويستقبل الماء البارد الكمية نفسها:
\[
m_2c(T_{eq}-T_2)=0{,}300\times4{,}18\times10^3\times26
\approx3{,}26\times10^4\,\text{J}.
\]
\textbf{4.} لا. للمادة نفسها يضرب $c$ طرفي موازنة الطاقة فيختزل. لا تعتمد درجة التوازن على $c$، بل هي متوسط درجات البداية موزوناً بالكتل. يتطلب قياس الحرارة النوعية بطريقة المزج مادة ذات سعة معلومة، كما في التمرين التالي.
\end{solution}
\end{exo}

\begin{exo}[تحديد الحرارة النوعية لمعدن بطريقة المزج]{calorimetrie-methode-melanges}
\seoready{true}
\keywords{المسعرية, طريقة المزج, الحرارة النوعية, المسعر, المكافئ المائي}

كما فعل جوزيف بلاك في القرن الثامن عشر، نريد قياس الحرارة النوعية $c_m$ لمعدن مجهول \emph{بطريقة المزج}. تُسخن عينة كتلتها $m=0{,}150\,\text{kg}$ إلى $T_m=95\,^\circ\text{C}$ ثم تُغمر سريعاً في مسعر يحوي $m_e=0{,}250\,\text{kg}$ من الماء، حيث $c_e=4{,}18\times10^3\,\text{J}\!\cdot\!\text{kg}^{-1}\!\cdot\!\text{K}^{-1}$ و$T_e=18\,^\circ\text{C}$. درجة التوازن المقاسة $T_{eq}=22\,^\circ\text{C}$. نهمل أولاً سعة الوعاء والمحراك والمحرار.
\begin{enumerate}
  \item اكتب حفظ الطاقة للنظام المعزول المؤلف من المعدن والماء، مع $c_m$ مجهولاً وحيداً.
  \item استنتج $c_m$ واحسبه. أي معدن شائع يقاربه؟ القيم المرجعية بوحدة $\text{J}\!\cdot\!\text{kg}^{-1}\!\cdot\!\text{K}^{-1}$: الألومنيوم $9{,}0\times10^2$، الحديد $4{,}5\times10^2$، النحاس $3{,}9\times10^2$، الرصاص $1{,}3\times10^2$.
  \item يمتص الوعاء بعض الحرارة. نمثل ذلك \emph{بمكافئ مائي} $\mu=1{,}5\times10^{-2}\,\text{kg}$، أي يتصرف المسعر ككتلة ماء إضافية $\mu$. أعد حساب $c_m$ وفسر اتجاه التصحيح.
  \item لماذا يجب غمر العينة سريعاً وعزل المسعر جيداً عن الهواء المحيط؟
\end{enumerate}
\begin{indication}
يفقد المعدن $Q_m=mc_m(T_m-T_{eq})$، ويستقبله الماء، وفي السؤال 3 المسعر أيضاً: $Q_m=(m_e+\mu)c_e(T_{eq}-T_e)$.
\end{indication}
\begin{solution}
\textbf{1.}
\[
mc_m(T_m-T_{eq})=m_ec_e(T_{eq}-T_e).
\]
\textbf{2.}
\[
c_m=\frac{m_ec_e(T_{eq}-T_e)}{m(T_m-T_{eq})}
=\frac{0{,}250\times4{,}18\times10^3\times(22-18)}{0{,}150\times(95-22)}
\approx3{,}82\times10^2\,\text{J}\!\cdot\!\text{kg}^{-1}\!\cdot\!\text{K}^{-1}.
\]
هذه القيمة قريبة جداً من قيمة النحاس.

\textbf{3.}
\[
c_m=\frac{(m_e+\mu)c_e(T_{eq}-T_e)}{m(T_m-T_{eq})}
=\frac{(0{,}250+0{,}015)\times4{,}18\times10^3\times4}{0{,}150\times73}
\approx4{,}05\times10^2\,\text{J}\!\cdot\!\text{kg}^{-1}\!\cdot\!\text{K}^{-1}.
\]
يزيد التصحيح $c_m$: إهمال الوعاء قلل تقدير الحرارة التي فقدها المعدن فعلاً، ومن ثم حرارته النوعية.

\textbf{4.} تفترض الطريقة نظاماً معزولاً طوال القياس. يقلل الغمر السريع تبادل العينة الساخنة للحرارة مع الهواء قبل بلوغ الماء، ويقلل العزل الجيد الفقد أثناء بلوغ التوازن. يشوه أي انتقال غير محسوب الموازنة وقيمة $c_m$. وهذا بالضبط ما راعاه بلاك، ثم لافوازييه ولابلاس في مسعرهما الجليدي.
\end{solution}
\end{exo}

\begin{exo}[السعرات الغذائية والقدرة الاستقلابية]{calorimetrie-calories-alimentaires}
\seoready{true}
\keywords{سعرة حرارية, كيلو سعرة, سعرة غذائية, المكافئ الميكانيكي, القدرة الاستقلابية, الوحدات}

يشير الملصق الغذائي إلى أن البالغ يستهلك في المتوسط نحو $2000\,\text{kcal}$ يومياً. الكيلو سعرة وحدة غير دولية ما زالت مستخدمة في التغذية؛ $1\,\text{kcal}=4{,}18\times10^3\,\text{J}$.
\begin{enumerate}
  \item حوّل الاستهلاك اليومي إلى جول.
  \item استنتج القدرة الاستقلابية المتوسطة بالواط، بافتراض استخدام الطاقة بانتظام خلال 24 ساعة.
  \item كُتب على لوح طاقة «$210\,\text{kcal}$». ما كتلة الماء عند $20\,^\circ\text{C}$ التي يمكن إيصالها إلى الغليان عند $100\,^\circ\text{C}$ بالطاقة نفسها إذا تحولت كلها إلى حرارة؟ خذ $c_{\text{water}}=4{,}18\times10^3\,\text{J}\!\cdot\!\text{kg}^{-1}\!\cdot\!\text{K}^{-1}$.
  \item في أي صور يستخدم جسم الإنسان هذه الطاقة فعلياً؟ أجب نوعياً.
\end{enumerate}
\begin{indication}
في السؤال 2 استخدم $\mathcal P=E/\Delta t$. في السؤال 3 حوّل أولاً إلى الجول، ثم استخرج $m$ من $Q=mc\Delta T$.
\end{indication}
\begin{solution}
\textbf{1.}
\[
E=2000\times4{,}18\times10^3=8{,}36\times10^6\,\text{J}.
\]
\textbf{2.}
\[
\mathcal P=\frac{8{,}36\times10^6}{8{,}64\times10^4}\approx97\,\text{W}.
\]
على مدار اليوم يستخدم الإنسان قدرة متوسطة تقارب قدرة مصباح متوهج. توضح رتبة المقدار المفاجئة غالباً فائدة تحويل السعرات الغذائية إلى وحدات فيزيائية مألوفة.

\textbf{3.}
\[
Q=210\times4{,}18\times10^3\approx8{,}78\times10^5\,\text{J}.
\]
مع $\Delta T=80\,\text{K}$،
\[
m=\frac{Q}{c\Delta T}
=\frac{8{,}78\times10^5}{4{,}18\times10^3\times80}
\approx2{,}63\,\text{kg}.
\]
إذن يحتوي لوح واحد، من حيث رتبة المقدار، على طاقة تكفي لغلي نحو $2{,}5\,\text{kg}$ من ماء الصنبور.

\textbf{4.} لا «يستهلك» الجسم الطاقة بمعنى إفنائها، بل يحول الطاقة الكيميائية للمغذيات. يُستخدم جزء في عمل العضلات والأعضاء وحفظ التدرجات الكهربائية للخلايا ونقل الجزيئات والتركيب الكيميائي، وقد يخزن جزء كغليكوجين أو دهون. في النهاية تتبدد معظم الطاقة المستخدمة حرارةً، فتسهم في حفظ حرارة الجسم قبل انتقالها إلى البيئة. ويخرج جزء صغير أيضاً طاقةً كيميائية باقية في الفضلات.
\end{solution}
\end{exo}
